% ============================================================
%  CHAPTER 6: CONCLUSION AND FUTURE WORK
% ============================================================
\chapter{Conclusion and Future Work}

This project proposed, implemented, and validated a Supply Chain Digital Twin for
the city of Nagpur, designed to address the problem of dynamic multi-objective
fleet routing under real-time constraints. The system combined a discrete-event
simulation engine with a three-tier AI architecture consisting of a Central PPO
Brain for macro-level strategy selection, per-truck Edge Q-Learning Agents for
emergency micro-rerouting, and an online Ridge Regression ETA Forecaster for
spatial-temporal travel time prediction.

The 365-day benchmark validated the system against a heuristic A* baseline under
identical conditions. The AI-enabled system fulfilled 323 additional orders,
delivered 474,708 more kilograms of cargo, and reduced retailer stockout events by
52.9\%, from 70,898 to 33,401. These improvements were achieved without any
degradation in cargo freshness or fuel efficiency, which were maintained at 99.7\%
and 78.7\% respectively in both configurations. Additionally, road accidents
reduced by 135 events as an emergent consequence of the ETA Forecaster learning
to route trucks away from high-accident zones during peak hours.

The reward normalization framework, which established a 5:1 freshness-to-fuel
penalty ratio to equalize the physical magnitudes of the two competing objectives,
was critical to enabling the PPO agent to learn balanced, context-sensitive
strategies. Without this normalization, the agent defaulted entirely to
fuel-conservative behavior and failed to exploit RSL-Priority and Speed-Priority
routing opportunities. The Anti-Thrashing GPS Lock was equally critical: before
this fix was applied, the fleet delivery rate collapsed to 49.8\% due to
continuous route recalculation. These two engineering contributions are not merely
implementation details; they represent fundamental design decisions that determine
whether the RL approach succeeds or fails in a real logistics deployment.

The results demonstrate that Pareto optimality is achievable in multi-objective
supply chain routing through reinforcement learning. An intelligent agent can
simultaneously improve throughput, supply continuity, and road safety without
requiring any explicit trade-offs between these objectives, provided that the
reward function is carefully designed and normalized to reflect the physical
realities of the operational environment.

% ================================================================
\section{Future Work}

Several extensions and improvements to the current system are identified for
future work.

\begin{enumerate}
    \item \textbf{Real-Time IoT Data Integration:} The current simulation relies on
    synthetic traffic and weather data. Connecting the Digital Twin to live IoT
    sensor streams from Nagpur's traffic management infrastructure, including
    inductive loop detectors and weather stations, would allow the ETA Forecaster
    to train on real-world data and produce more accurate travel time predictions.

    \item \textbf{Multi-Seed Statistical Validation:} The 365-day benchmark was
    conducted on a single fixed random seed (seed=42) for reproducibility. Running
    the benchmark across a set of 5 to 10 different seeds would produce mean and
    standard deviation estimates for each performance metric, providing a statistically
    rigorous basis for the claims made in this thesis.

    \item \textbf{Testing Diverse Perishables:} The current simulation focuses
    exclusively on a single perishable product profile (oranges) with a fixed
    decay rate. Future iterations should test the AI against a diverse mix of
    goods, such as pharmaceuticals, dairy, and frozen foods, each requiring
    distinct temperature profiles and varying Remaining Shelf Life (RSL)
    sensitivities.

    \item \textbf{Intercity Travel and Cold Storage Logistics:} Extending the
    Nagpur city graph to include intercity highway networks and regional cold
    storage hubs. This would allow the reinforcement learning architecture to
    optimize long-haul logistics and manage the energy consumption of
    refrigerated containers over multi-day journeys.

    \item \textbf{Deep RL for the Edge Agent:} The Edge Agent currently uses a
    small tabular Q-table that requires manual discretization of the local hazard
    state space. Replacing this with a Deep Q-Network would allow the Edge Agent
    to generalize across new intersection configurations and hazard types that were
    not encountered during training, improving its adaptability in deployment.
\end{enumerate}
